% =====================================================================
% main.tex
% Responsabilidad única: ensamblar el documento completo a partir de los
% módulos (preámbulo, portada y las diez secciones de ejercicios). No
% contiene desarrollo matemático ni resultados propios — solo estructura.
%
% Para compilar:  latexmk -pdf main.tex
% =====================================================================
\documentclass[11pt]{article}

% =====================================================================
% preambulo.tex
% Responsabilidad única: cargar paquetes y definir configuración/estilos
% reutilizables (colores de enlaces, entornos, listas). No contiene
% contenido propio del taller — si algo del "look" del documento debe
% cambiar, este es el único archivo que hay que tocar.
% =====================================================================

\usepackage[utf8]{inputenc}
\usepackage[T1]{fontenc}
\usepackage[spanish,es-noshorthands,es-tabla]{babel}
\usepackage[a4paper,margin=2.5cm]{geometry}

% --- Matemáticas ------------------------------------------------------
\usepackage{amsmath}
\usepackage{amssymb}
\usepackage{amsthm}
\usepackage{mathtools}

% --- Tablas y figuras ---------------------------------------------------
\usepackage{graphicx}
\usepackage{booktabs}
\usepackage{array}
\usepackage{multirow}
\usepackage{caption}
\usepackage{float}
\usepackage{pgfplots}
\pgfplotsset{compat=1.18}

% --- Listas con formato a., b., c. ------------------------------------
\usepackage{enumitem}

% --- Diagramas vectoriales -----------------------------------------------
\usepackage{tikz}
\usetikzlibrary{arrows.meta, positioning, calc}
% Misma paleta usada en las figuras generadas con matplotlib (Ejercicios
% 6-10), para que los diagramas TikZ se vean como parte del mismo documento.
\definecolor{classblue}{HTML}{4C72B0}
\definecolor{classred}{HTML}{C44E52}
\definecolor{classgreen}{HTML}{55A868}

% --- Color y enlaces ----------------------------------------------------
\usepackage{xcolor}
\definecolor{urosariored}{RGB}{255,11,38} % rojo exacto muestreado del escudo
\usepackage{hyperref}
\hypersetup{
    colorlinks = true,
    linkcolor  = urosariored!75!black,
    citecolor  = urosariored!75!black,
    urlcolor   = urosariored!75!black,
    pdftitle   = {Taller 2 -- Modelos de aprendizaje supervisado},
    pdfauthor  = {Tania Fernanda Mora Cubillos, Over Haider Castrillon Valencia},
}

% --- Espaciado de tablas ------------------------------------------------
\renewcommand{\arraystretch}{1.25}

% --- Pie de página institucional -----------------------------------------
% Responsabilidad de este bloque: mostrar el escudo de la Universidad del
% Rosario, pequeño, en la esquina inferior derecha de cada página del
% cuerpo del documento (la portada define su propio pie por separado en
% portada.tex, ya que \begin{titlepage} usa página sin estilo).
\usepackage{fancyhdr}
\pagestyle{fancy}
\fancyhf{}
\renewcommand{\headrulewidth}{0pt}
\renewcommand{\footrulewidth}{0pt}
\fancyfoot[C]{\thepage}
\fancyfoot[R]{\includegraphics[height=0.42cm]{figuras/logo_urosario_icono.png}}

% --- Encabezados de subsección sin numerar, para incisos a-e -----------
\setlist[enumerate,1]{label=\alph*., leftmargin=1.6em}

% --- Cifras generadas por el notebook -------------------------------------
% Responsabilidad de este bloque: cargar las macros con los resultados
% numéricos de los ejercicios 6-10 que escribe el notebook al ejecutarse
% (generado/resultados.tex), de modo que el texto cite \macro en lugar de
% copiar números a mano. Si el archivo aún no existe, se compila igual.
\InputIfFileExists{generado/resultados}{}{}


\begin{document}

% =====================================================================
% portada.tex
% Responsabilidad única: presentar los datos identificativos del taller
% (título, curso, estudiantes, profesor, fecha). No contiene desarrollo
% matemático ni resultados.
% =====================================================================

\begin{titlepage}
    \centering
    \vspace*{2cm}

    {\LARGE \textbf{Taller 2 --- Modelos de aprendizaje supervisado}}\\[0.4cm]
    {\normalsize\itshape Regresión lineal simple y múltiple, regresión logística y árboles de decisión}\\[0.6cm]
    {\large Aprendizaje Automático de Máquina I}\\[0.15cm]
    {\large Semestre 2026-2}\\[1.6cm]

    \rule{0.8\textwidth}{0.4pt}\\[0.9cm]

    \begin{tabular}{rl}
        \textbf{Estudiantes:}   & Tania Fernanda Mora Cubillos \\
                                & \small\texttt{taniaf.mora@urosario.edu.co} \\[0.35em]
                                & Over Haider Castrillón Valencia \\
                                & \small\texttt{over.castrillon@urosario.edu.co} \\[0.35em]
        \textbf{Programa:}      & Maestría en Matemáticas y Ciencias de la Computación \\[0.35em]
        \textbf{Profesor:}      & Daniel Felipe Rambaut Lemus \\[0.35em]
        \textbf{Modalidad:}     & Grupal \\[0.35em]
        \textbf{Fecha:}         & \today \\
    \end{tabular}

    \vfill
    \includegraphics[width=2.6cm]{figuras/logo_urosario.png}
\end{titlepage}


\tableofcontents
\newpage

% ---------------------------------------------------------------------
% Parte I — Componente teórico-matemático
% ---------------------------------------------------------------------
\section{Componente teórico-matemático}
% =====================================================================
% ejercicio01.tex
% Responsabilidad única: derivar las ecuaciones normales de la recta de mínimos cuadrados, ajustar
% los datos didácticos, verificar propiedades de los residuos y estudiar el
% efecto de una observación influyente.
% =====================================================================
\subsection{Ejercicio 1 --- Regresión lineal simple: ajuste e influencia de una observación}

\begin{enumerate}

\item \textbf{Ecuaciones normales y estimadores.}
% TODO

\item \textbf{Coeficientes, predicciones y residuos.}
% TODO

\item \textbf{MSE, $R^2$ y comparación con $\bar y$.}
% TODO

\item \textbf{Observación influyente ($y_5 = 12$).}
% TODO

\end{enumerate}

% =====================================================================
% ejercicio02.tex
% Responsabilidad única: plantear y resolver el sistema normal
% X^T X beta = X^T y, interpretar coeficientes y discutir colinealidad
% perfecta y el efecto de añadir variables sobre el error de
% entrenamiento y de generalización. Cifras verificadas con
% verificacion/teoria.py.
% =====================================================================
\subsection{Ejercicio 2 --- Regresión lineal múltiple: solución matricial y colinealidad}

Se modela la demanda $y$ (centenas de unidades) como
$\hat y=\beta_0+\beta_1x_1+\beta_2x_2$, donde $x_1$ es la inversión
publicitaria (unidades de presupuesto) y $x_2\in\{0,1\}$ indica si hubo
promoción.

\begin{enumerate}

\item \textbf{Matriz $X$, $X^\top X$ y $X^\top y$.}
Agregando la columna de unos del intercepto,
\[
    X = \begin{pmatrix}
        1 & 0 & 0\\ 1 & 1 & 0\\ 1 & 2 & 0\\ 1 & 0 & 1\\ 1 & 1 & 1\\ 1 & 2 & 1
    \end{pmatrix}\in\mathbb{R}^{6\times 3},
    \qquad
    y = \begin{pmatrix} 2\\4\\5\\5\\6\\8 \end{pmatrix}\in\mathbb{R}^{6\times 1},
    \qquad
    \beta\in\mathbb{R}^{3\times1},
\]
y cada entrada de $X^\top X$ es un producto punto entre columnas de $X$
(por ejemplo $\sum_i x_{i1}^2=0+1+4+0+1+4=10$ y
$\sum_i x_{i1}x_{i2}=0+1+2=3$), así como cada entrada de $X^\top y$ es
el producto punto de una columna con $y$, de forma que
\[
    X^\top X = \begin{pmatrix}
        n & \sum x_{1} & \sum x_{2}\\
        \sum x_{1} & \sum x_{1}^2 & \sum x_{1}x_{2}\\
        \sum x_{2} & \sum x_{1}x_{2} & \sum x_{2}^2
    \end{pmatrix}
    = \begin{pmatrix} 6 & 6 & 3\\ 6 & 10 & 3\\ 3 & 3 & 3 \end{pmatrix},
    \qquad
    X^\top y = \begin{pmatrix} \sum y\\ \sum x_1y\\ \sum x_2y \end{pmatrix}
    = \begin{pmatrix} 30\\ 36\\ 19 \end{pmatrix}.
\]

\item \textbf{Derivación y solución del sistema normal.}
Expandiendo la función objetivo,
\[
    S(\beta) = (y-X\beta)^\top(y-X\beta)
    = y^\top y - 2\beta^\top X^\top y + \beta^\top X^\top X\beta ,
\]
y usando las reglas $\nabla_\beta(\beta^\top a)=a$ y
$\nabla_\beta(\beta^\top A\beta)=2A\beta$ para $A$ simétrica, se obtiene
\[
    \nabla_\beta S = -2X^\top y + 2X^\top X\beta = 0
    \;\Longrightarrow\;
    X^\top X\hat\beta = X^\top y ,
\]
que corresponde a un mínimo porque la hessiana $2X^\top X$ es semidefinida
positiva ($v^\top X^\top Xv=\lVert Xv\rVert^2\ge0$), y definida positiva
cuando las columnas de $X$ son linealmente independientes, como en este
caso ($\det X^\top X=36\neq0$). El sistema a resolver es
\[
    \begin{aligned}
        6\beta_0 + 6\beta_1 + 3\beta_2 &= 30 &\quad&(\mathrm{E}_1)\\
        6\beta_0 + 10\beta_1 + 3\beta_2 &= 36 &&(\mathrm{E}_2)\\
        3\beta_0 + 3\beta_1 + 3\beta_2 &= 19 &&(\mathrm{E}_3)
    \end{aligned}
\]
y se procede por eliminación: primeramente $\mathrm{E}_2-\mathrm{E}_1$
da $4\beta_1=6$, luego $\mathrm{E}_1-2\mathrm{E}_3$ da $-3\beta_2=-8$, y
finalmente al sustituir ambos en $\mathrm{E}_3$ se despeja $\beta_0$:
\[
    \hat\beta_1 = \frac{3}{2} = 1.5, \qquad
    \hat\beta_2 = \frac{8}{3} \approx 2.6667, \qquad
    3\hat\beta_0 = 19 - 4.5 - 8 = 6.5
    \;\Rightarrow\; \hat\beta_0 = \frac{13}{6} \approx 2.1667 .
\]

\item \textbf{Interpretación, predicción y MSE.}
\begin{itemize}
    \item $\hat\beta_1=1.5$: manteniendo fija la condición de promoción,
    cada unidad adicional de presupuesto publicitario se asocia con
    $1.5$ centenas ($150$ unidades) más de demanda predicha.
    \item $\hat\beta_2\approx2.6667$: para un mismo nivel de inversión
    publicitaria, realizar la promoción se asocia con unas $2.6667$
    centenas (aproximadamente $267$ unidades) más de demanda predicha.
\end{itemize}
Para $(x_1,x_2)=(1.5,1)$ se predice
$\hat y=\tfrac{13}{6}+1.5\cdot1.5+\tfrac83=\tfrac{85}{12}\approx7.0833$
centenas, es decir, unas $708$ unidades. Con las predicciones de
entrenamiento
\[
    \hat y\approx(2.1667,\;3.6667,\;5.1667,\;4.8333,\;6.3333,\;7.8333),
\]
los residuos son $e=\pm\tfrac16$ o $\pm\tfrac13$ y
\[
    MSE = \frac{1}{6}\Bigl(4\cdot\tfrac{1}{36} + 2\cdot\tfrac{1}{9}\Bigr)
    = \frac{1}{6}\cdot\frac{1}{3} = \frac{1}{18} \approx 0.0556 .
\]
Estos coeficientes no prueban un efecto causal, porque solo describen
la mejor combinación lineal \emph{en estos datos}, y los datos no
provienen (hasta donde se sabe) de un experimento donde la inversión o
la promoción se asignaran al azar, de forma que podría existir una
variable omitida (temporada, precio, competencia) que mueva a la vez la
decisión de invertir y la demanda; además con $n=6$ observaciones
didácticas no hay forma de separar el efecto real del ruido, así que la
lectura correcta es de \emph{asociación condicional}, no de
intervención.

\item \textbf{Colinealidad perfecta ($x_3 = 2x_1$).}
Si se agrega $x_3=2x_1$, la nueva matriz
$\tilde X=[\mathbf{1},x_1,x_2,x_3]$ satisface
$0\cdot\mathbf{1}+2\,x_1+0\cdot x_2-1\cdot x_3=\mathbf{0}$, es decir,
existe $v=(0,2,0,-1)^\top\neq0$ con $\tilde Xv=0$, por lo que sus
columnas son linealmente dependientes, $\operatorname{rango}(\tilde X)=3<4$
y $\tilde X^\top\tilde X$ (de $4\times4$) es singular. En consecuencia,
si $\hat\beta$ es una solución, también lo es $\hat\beta+c\,v$ para todo
$c\in\mathbb{R}$, porque $\tilde X(\hat\beta+cv)=\tilde X\hat\beta$:
cualquier combinación con $\beta_1+2\beta_3=1.5$ (y los mismos
$\beta_0,\beta_2$) produce \textbf{exactamente las mismas
predicciones}, de modo que el ajuste sigue siendo único en términos de
$\hat y$, pero no en términos de coeficientes, que dejan de ser
interpretables individualmente (no tiene sentido preguntar ``el efecto
de $x_1$ con $x_3$ fijo'' si $x_3$ no puede quedarse fijo mientras $x_1$
cambia). Para ajustar sin invertir una matriz singular se puede, en
orden de preferencia:
\begin{itemize}
    \item Eliminar la variable redundante: Se descarta $x_3$ (o $x_1$),
    lo que devuelve exactamente el modelo del inciso b.
    \item Usar la pseudoinversa de Moore--Penrose: Se calcula
    $\hat\beta=\tilde X^{+}y$ (mediante descomposición en valores
    singulares, \texttt{np.linalg.lstsq} o \texttt{pinv}), que elige
    entre las infinitas soluciones la de norma mínima, en este caso
    $\hat\beta\approx(2.1667,\,0.3,\,2.6667,\,0.6)$, con
    $0.3+2\cdot0.6=1.5$ y las mismas predicciones.
    \item Regularizar: Se resuelve
    $(\tilde X^\top\tilde X+\lambda I)\beta=\tilde X^\top y$ con
    $\lambda>0$ (regresión ridge), sistema que siempre es invertible.
\end{itemize}

\item \textbf{Más variables: error de entrenamiento vs.\ datos nuevos.}
Si se añade una variable $x_{p+1}$, el modelo ampliado contiene al
original como caso particular (basta fijar $\beta_{p+1}=0$), de forma
que el mínimo del SSE sobre un conjunto más grande de candidatos no
puede ser mayor:
\[
    \min_{\beta\in\mathbb{R}^{p+2}} \lVert y-\tilde X\beta\rVert^2
    \;\le\;
    \lVert y-\tilde X(\hat\beta_{\text{orig}},0)\rVert^2
    = \min_{\beta\in\mathbb{R}^{p+1}} \lVert y-X\beta\rVert^2 ,
\]
y por eso el error de entrenamiento nunca aumenta (en el peor caso se
mantiene igual, cuando la nueva variable no aporta nada en la muestra).
Sin embargo, en datos nuevos sí puede empeorar, porque el nuevo
coeficiente también se ajusta al ruido particular de la muestra de
entrenamiento, lo cual aumenta la varianza de las predicciones sin
reducir necesariamente el sesgo (sobreajuste), y con pocas
observaciones o predictores correlacionados este efecto se vuelve más
fuerte, razón por la cual en la Parte II la comparación entre modelos se
decide con un conjunto de validación y no con el error de
entrenamiento.

\end{enumerate}

% =====================================================================
% ejercicio03.tex
% Responsabilidad única: relacionar la sigmoide con los log-odds,
% clasificar tres máquinas bajo distintos umbrales, deducir la frontera
% de decisión y el umbral que minimiza el costo esperado. Cifras
% verificadas con verificacion/teoria.py.
% =====================================================================
\subsection{Ejercicio 3 --- Regresión logística: probabilidad, odds y umbral}

El modelo estima la probabilidad de falla en la próxima semana ($Y=1$)
a partir de los años de uso $x_1$ y las alertas del último mes $x_2$:
\[
    z = -3 + 0.6x_1 + 1.2x_2, \qquad p = \sigma(z) = \frac{1}{1+e^{-z}} .
\]

\begin{enumerate}

\item \textbf{Log-odds y factor de odds.}
A partir de la definición de $p$,
\[
    1-p = 1 - \frac{1}{1+e^{-z}} = \frac{e^{-z}}{1+e^{-z}}
    \;\Longrightarrow\;
    \frac{p}{1-p} = \frac{1}{e^{-z}} = e^{z}
    \;\Longrightarrow\;
    \log\frac{p}{1-p} = z ,
\]
es decir, el modelo es \emph{lineal en los log-odds}, no en la
probabilidad. Si $x_1$ aumenta en una unidad con $x_2$ fijo, $z$ aumenta
exactamente $0.6$ y los odds pasan de $e^{z}$ a $e^{z+0.6}$, de forma que
se multiplican por
\[
    e^{0.6} \approx 1.8221 ,
\]
lo cual equivale a un aumento del 82.21\,\% \emph{en los odds},
constante sin importar el punto de partida. Esto no debe confundirse con
un cambio en puntos porcentuales de la probabilidad, que sí depende del
punto de partida porque la sigmoide no es lineal: pasar de $(2,1)$ a
$(3,1)$ lleva $p$ de $0.3543$ a $0.5000$ ($+14.57$ puntos porcentuales),
mientras que pasar de $(4,2)$ a $(5,2)$ lleva $p$ de $0.8581$ a $0.9168$
(apenas $+5.87$ puntos), aun cuando en ambos casos los odds se
multiplican por el mismo $1.8221$.

\item \textbf{Probabilidades y clases para $A$, $B$, $C$.}
Usando $\hat y=1$ cuando $p\ge t$:

\begin{table}[H]
\centering
\begin{tabular}{crrrrcc}
\toprule
Máquina & $(x_1,x_2)$ & $z$ & $p=\sigma(z)$ & odds & $\hat y$ ($t=0.5$) & $\hat y$ ($t=0.8$) \\
\midrule
$A$ & $(2,1)$ & $-3+1.2+1.2=-0.6$ & 0.3543 & 0.5488 & 0 & 0 \\
$B$ & $(5,0)$ & $-3+3.0+0\phantom{.0}=\phantom{-}0.0$ & 0.5000 & 1.0000 & 1 & 0 \\
$C$ & $(4,2)$ & $-3+2.4+2.4=\phantom{-}1.8$ & 0.8581 & 6.0496 & 1 & 1 \\
\bottomrule
\end{tabular}
\caption{Puntaje lineal, probabilidad de falla y clase predicha por umbral.}
\label{tab:ej3-maquinas}
\end{table}

Cabe mencionar el caso de $B$, que cae exactamente sobre la frontera de
$t=0.5$ ($p=0.5$): con la regla $p\ge t$ se clasifica como falla, pero
con una regla estricta $p>t$ se clasificaría como no falla, por lo cual
la convención de desempate debe declararse explícitamente. Al subir el
umbral a $0.8$ solo la máquina $C$ conserva la predicción de falla.

\item \textbf{Frontera de decisión para un umbral general.}
Como $\sigma$ es estrictamente creciente, $p\ge t$ equivale a
$z\ge\sigma^{-1}(t)=\log\frac{t}{1-t}$, por lo que la frontera es la
recta
\[
    0.6x_1 + 1.2x_2 = 3 + \log\frac{t}{1-t}
    \quad\Longleftrightarrow\quad
    x_2 = \frac{3+\log\frac{t}{1-t}}{1.2} - 0.5\,x_1 ,
\]
y se predice falla en el semiplano por encima de ella. Particularizando,
\[
    \begin{aligned}
    t=0.5:&\quad 0.6x_1+1.2x_2 = 3 &&\Longleftrightarrow\quad x_2 = 2.5 - 0.5x_1,\\
    t=0.8:&\quad 0.6x_1+1.2x_2 = 3+\log 4 \approx 4.3863 &&\Longleftrightarrow\quad x_2 \approx 3.6552 - 0.5x_1.
    \end{aligned}
\]
Al modificar solamente el umbral \textbf{la frontera se desplaza
paralelamente} (cambia el término independiente, de forma que con
$t=0.8$ se exige más evidencia y la región de falla se reduce), pero su
pendiente $-0.5=-\beta_1/\beta_2$ y su orientación quedan iguales, así
como \textbf{las probabilidades no cambian en absoluto}: $p(x)$ depende
solo de $\beta$, de modo que el umbral es una decisión posterior sobre
cómo convertir probabilidades en clases, no una propiedad del modelo
ajustado.

\begin{figure}[H]
\centering
\begin{tikzpicture}
\begin{axis}[
    width=0.66\textwidth, height=6cm,
    xmin=0, xmax=7.5, ymin=0, ymax=4.2,
    xlabel={$x_1$ (años de uso)}, ylabel={$x_2$ (alertas)},
    legend pos=north east, legend cell align=left,
    legend style={font=\small}, tick label style={font=\small},
]
\addplot[domain=0:7.5, thick, color=classblue] {2.5 - 0.5*x};
\addlegendentry{$t=0.5$}
\addplot[domain=0:7.5, thick, dashed, color=classred] {3.6552 - 0.5*x};
\addlegendentry{$t=0.8$}
\addplot[domain=0:7.5, thick, dotted, color=classgreen] {1.3448 - 0.5*x};
\addlegendentry{$t=0.2$ (inciso d)}
\addplot[only marks, mark=*, mark size=2.4pt, color=black,
         nodes near coords, point meta=explicit symbolic,
         every node near coord/.append style={font=\small, anchor=south west}]
    coordinates {(2,1)[$A$] (5,0)[$B$] (4,2)[$C$]};
\end{axis}
\end{tikzpicture}
\caption{Fronteras de decisión $p=t$ en el plano $(x_1,x_2)$: rectas
paralelas de pendiente $-0.5$; por encima de cada recta se predice falla.}
\end{figure}

\item \textbf{Umbral de costo mínimo.}
Con probabilidades calibradas, si se predice $0$ el único error posible
es un falso negativo (con probabilidad $p$ y costo $8$), y si se predice
$1$ el único error posible es un falso positivo (con probabilidad $1-p$
y costo $2$), de forma que
\[
    \mathbb{E}[\text{costo}\mid\hat y=0] = 8p,
    \qquad
    \mathbb{E}[\text{costo}\mid\hat y=1] = 2(1-p),
\]
y conviene predecir $1$ cuando $2(1-p)\le8p$, es decir, cuando
$p\ge 2/(2+8)=\mathbf{0.2}$ (en general,
$t^*=c_{FP}/(c_{FP}+c_{FN})$). Como un falso negativo cuesta
cuatro veces más que un falso positivo, el umbral óptimo baja de $0.5$
a $0.2$, y aplicándolo, $A$ ($p=0.3543$), $B$ ($p=0.5$) y $C$
($p=0.8581$) superan $0.2$, por lo cual \textbf{las tres máquinas se
clasifican como falla}, incluida $A$, que con $t=0.5$ no se marcaba
(su frontera aparece punteada en la figura anterior,
$x_2\approx1.3448-0.5x_1$).

\item \textbf{Mínimos cuadrados vs.\ sigmoide en $[0,1]$.}
Un modelo $\hat y=\beta_0+\beta^\top x$ es una función afín no acotada:
si algún $\beta_j\neq0$, al crecer o decrecer $x_j$ lo suficiente
$\hat y$ supera $1$ o cae por debajo de $0$, y mínimos cuadrados no
incluye ninguna restricción que lo evite, pues solo minimiza la
distancia cuadrática a etiquetas $0/1$ (por ejemplo, ajustando MCO a
$x=(0,1,2,3,10)$ con $y=(0,0,1,1,1)$ se obtiene
$\hat y\approx0.3248+0.0860x$, que predice $\hat y(15)\approx1.6146$, un
valor imposible como probabilidad). En cambio, la sigmoide cumple
$\sigma(z)=1/(1+e^{-z})\in(0,1)$ para todo $z\in\mathbb{R}$, porque
$e^{-z}>0$ hace que el denominador sea siempre mayor que $1$, de modo
que $\sigma$ transforma el puntaje lineal no acotado en una
probabilidad válida, tendiendo a $0$ y a $1$ solo asintóticamente.

\end{enumerate}

% =====================================================================
% ejercicio04.tex
% Responsabilidad única: evaluar la entropía cruzada, deducir su gradiente, ejecutar un paso de
% descenso de gradiente y analizar la penalización L2 y el parámetro $C$.
% =====================================================================
\subsection{Ejercicio 4 --- Regresión logística: pérdida, gradiente y regularización}

\begin{enumerate}

\item \textbf{Probabilidades y pérdida inicial.}
% TODO

\item \textbf{Gradiente de $J$.}
% TODO

\item \textbf{Un paso de descenso de gradiente.}
% TODO

\item \textbf{Penalización $\lambda\beta_1^2/2$.}
% TODO

\item \textbf{$C$, regularización y escala de entrada.}
% TODO

\end{enumerate}

% =====================================================================
% ejercicio05.tex
% Responsabilidad única: comparar divisiones por reducción de Gini, asignar clases en hojas,
% elegir una división de regresión por suma de errores cuadrados y
% discutir extrapolación y control de complejidad.
% =====================================================================
\subsection{Ejercicio 5 --- Árboles de decisión: división, predicción y complejidad}

\begin{enumerate}

\item \textbf{Impureza Gini y reducción ponderada.}
% TODO

\item \textbf{Clases y probabilidades en las hojas.}
% TODO

\item \textbf{División de regresión por SSE.}
% TODO

\item \textbf{La media minimiza el SSE; no extrapolación.}
% TODO

\item \textbf{\texttt{max\_depth}, \texttt{min\_samples\_leaf} y sobreajuste.}
% TODO

\end{enumerate}


\newpage
% ---------------------------------------------------------------------
% Parte II — Componente práctico-computacional
% ---------------------------------------------------------------------
\section{Componente práctico-computacional}
\label{sec:practico}

% TODO(rama práctica): párrafo de datos (load_diabetes con scaled=False,
% load_breast_cancer recodificado 1 = maligno) y protocolo 60/20/20 con
% random_state=42; tamaños de partición vía macros de generado/resultados.

% =====================================================================
% ejercicio06.tex
% Responsabilidad única: explorar el caso diabetes en entrenamiento, ajustar la regresión simple
% con \texttt{bmi} (NumPy y scikit-learn), compararla con el predictor
% constante y analizar residuos de validación.
% =====================================================================
\subsection{Ejercicio 6 --- Caso A: regresión lineal simple y modelo de referencia}

\begin{enumerate}

\item \textbf{Exploración en entrenamiento.}
% TODO

\item \textbf{Implementación con NumPy vs.\ \texttt{LinearRegression}.}
% TODO

\item \textbf{Comparación con \texttt{DummyRegressor}.}
% TODO

\item \textbf{Residuos de validación.}
% TODO

\item \textbf{Interpretación de coeficientes y causalidad.}
% TODO

\end{enumerate}

% =====================================================================
% ejercicio07.tex
% Responsabilidad única: ajustar la regresión múltiple estandarizada y la cuadrícula de árboles de
% regresión, y seleccionar con validación el candidato final de regresión.
% =====================================================================
\subsection{Ejercicio 7 --- Caso A: regresión múltiple y árbol de regresión}

\begin{enumerate}

\item \textbf{Regresión múltiple con \texttt{Pipeline}.}
% TODO

\item \textbf{Coeficientes estandarizados y correlaciones.}
% TODO

\item \textbf{Cuadrícula de \texttt{DecisionTreeRegressor}.}
% TODO

\item \textbf{Candidato final de regresión.}
% TODO

\item \textbf{No linealidad del árbol y brecha entrenamiento--validación.}
% TODO

\end{enumerate}

% =====================================================================
% ejercicio08.tex
% Responsabilidad única: ajustar la regresión logística regularizada en el caso cáncer de mama,
% seleccionar $C$ por ROC-AUC de validación y fijar el umbral por costo.
% =====================================================================
\subsection{Ejercicio 8 --- Caso B: regresión logística y elección del umbral}

\begin{enumerate}

\item \textbf{Particiones, recodificación y referencia.}
% TODO

\item \textbf{Selección de $C$.}
% TODO

\item \textbf{Umbral por costo $L = FP + 5FN$.}
% TODO

\item \textbf{Precisión y recall vs.\ umbral.}
% TODO

\item \textbf{Coeficientes estandarizados y odds.}
% TODO

\end{enumerate}

% =====================================================================
% ejercicio09.tex
% Responsabilidad única: ajustar la cuadrícula de árboles de clasificación, elegir árbol y umbral
% con validación, comparar con la logística e interpretar reglas e
% importancias.
% =====================================================================
\subsection{Ejercicio 9 --- Caso B: árbol de clasificación e interpretación}

\begin{enumerate}

\item \textbf{Cuadrícula de \texttt{DecisionTreeClassifier}.}
% TODO

\item \textbf{Selección del árbol.}
% TODO

\item \textbf{Umbral del árbol y candidato final de clasificación.}
% TODO

\item \textbf{Visualización y regla de una hoja.}
% TODO

\item \textbf{Importancias, coeficientes y sobreajuste.}
% TODO

\end{enumerate}

% =====================================================================
% ejercicio10.tex
% Responsabilidad única: registrar los candidatos congelados, evaluarlos una única vez en prueba
% y redactar la recomendación final.
% =====================================================================
\subsection{Ejercicio 10 --- Evaluación final y recomendación sustentada}

\begin{enumerate}

\item \textbf{Regresión en prueba.}
% TODO

\item \textbf{Clasificación en prueba.}
% TODO

\item \textbf{Matrices de confusión y costo.}
% TODO

\item \textbf{Recomendación.}
% TODO

\end{enumerate}


\end{document}
